\documentclass[uplatex,dvipdfmx,ja=standard,a4paper,11pt]{bxjsarticle}
\usepackage{geometry}
\geometry{margin=25mm}

\begin{document}

\begin{center}
    {\huge 生成AI利用申告書}
\end{center}

\vspace{1.5em}

\noindent
25G1110 早川 和輝 \\
2026年4月27日 \\
使用した AIツール：Gemini

\vspace{2em}

\noindent
今回の課題において，Arduinoのシリアルプロッタで正弦波を綺麗に表示させるために，作成したプログラム内のどの出力部分が邪魔になっているかの特定の確認に生成AIを使用した．

\vspace{1em}

\noindent
作成したArduinoのソースコード全体をプロンプトに入力した上で，「このプログラムで正弦波を表示するために邪魔な箇所はどこですか」という質問と指示を出した．

\vspace{1em}

\noindent
「一部修正して使った」に該当する．AIから提示された修正案のうち，シリアルプロッタで正弦波の波形だけをシンプルに確認するための出力方法を採用し，自身のプログラムの不要な文字出力部分を削除して書き換えたためである．

\vspace{1em}

\noindent
AIに指摘された不要な文字列出力の箇所を削除し，数値のみを出力するようにコードを修正した後，実際にArduinoボードにプログラムを書き込んで実行した．その後，Arduino IDEのシリアルプロッタを起動し，プログラムのボーレートと合致していることを確認した上で，意図した通りに綺麗な正弦波のグラフが描画されることを目視で確認した．

\vspace{1em}

\noindent
間違いは「なかった」と判断した．AIの指摘した通り，シリアルプロッタは文字列が混ざったり複数の不要なデータが送信されたりすると波形が正しく描画されないという仕様は事実であり，実際に不要なテキスト出力を削った結果，波形が正常に表示されたためである．

\vspace{1em}

\noindent
AIからは波形のみを表示するパターンと，最大値や最小値を同時に出力してグラフの縦軸のスケールを固定するパターンの2種類が提案された．今回はまず波形の形状をシンプルに確認することを最優先とし，単一の変数のみを出力する方式を自身で選択して実装した．

\vfill
\begin{center}
1
\end{center}

\end{document}